\chapter{Docking a ligand}
\label{ch:docking}

\term{Docking} predicts where a small molecule sits in a protein's pocket, and
how tightly. ProPrep's docking tool prepares both the protein and the molecule
for you and runs AutoDock Vina. This chapter follows a complete run on HIV
protease (\pdb{1HSG}), putting its own drug molecule back where the crystal
structure found it.

\begin{glance}
\glanceitem{What it does}{Puts a small molecule into a pocket of your protein
in many positions, scores each one, and shows you the best few.}
\glanceitem{You need it when}{You want to know how and how well a molecule
binds, or to rank many molecules against one protein.}
\glanceitem{Before you start}{Load a structure with \key{1a}. The tool needs
\texttt{vina} installed, and \texttt{autogrid4} as well for AD4 scoring.}
\glanceitem{Where to find it}{\menupath{5d}{Molecular Docking}, then \key{1}
\textsf{Dock a ligand into the structure}.}
\glanceitem{You get}{A folder \file{docking/docking_<date>_<time>/} with
\file{poses.sdf}, one file per pose, and \file{docking.json}.}
\end{glance}

\section{What docking can and cannot tell you}

Docking searches for positions and shapes of the molecule inside a box you
choose, and gives each one a score in kcal/mol. The score is a rough estimate
of binding strength, not a measurement.

\begin{itemize}
  \item \textbf{It can} suggest how a molecule is likely to sit in a pocket,
        which parts of the protein it touches, and which of several molecules
        looks more promising.
  \item \textbf{It cannot} give a reliable binding energy, nor tell you about
        anything it does not model: the protein barely moves, water is left
        out, and the molecule's charges are taken as given.
\end{itemize}

The honest test is \term{redocking}: take a structure solved with its molecule
in place, remove the molecule, and dock it back. The \term{RMSD} (root mean
square deviation, the average distance between matching atoms) then tells you
how close the prediction came to reality. Under about 2\,\AA{} means the
method reproduced the known answer. ProPrep shows this column
automatically when the molecule came from the structure.

\begin{behindscenes}
ProPrep prepares both sides with Meeko and pdb2pqr and writes them in
\term{PDBQT} format, the file type AutoDock expects. No force-field parameters
and no tLEaP are involved, so docking needs far less preparation than a
simulation.
\end{behindscenes}

\begin{needsoftware}
Docking needs AutoDock Vina (the program \texttt{vina}) on your computer, and
\texttt{autogrid4} as well if you want the \texttt{ad4} scoring function. Both
were present on our test machine, so all three scoring functions worked. The
molecules are prepared by parts of ProPrep's own installation (Meeko, RDKit
and pdb2pqr), so nothing else is needed. See
Chapter~\ref{ch:requirements} for how to check what you have.
\end{needsoftware}

\section{Step by step: redocking MK1 into HIV protease}
\label{sec:docking-steps}

\pdb{1HSG} is HIV protease with the drug indinavir (residue \texttt{MK1}) in
its active site. Start ProPrep with \texttt{proprep --full-menu} and load the
structure first: \key{1a}, \key{1}, \key{1}, \key{1}, type \typed{1HSG},
\key{1}, then \key{m}.

\begin{tip}
Open \key{5d} with no structure loaded and you only get
\prompt{Needs a loaded structure} (with a warning sign) and the menu again.
Load the structure
first.
\end{tip}

Type \key{5d}, then \key{1}. The dashboard lists nine setup steps and tells
you what is still missing:

\begin{screen}
  Current setup
  1 Receptor:     not set
  2 Ligand:       not set
  3 Kept HETATM:  not set
  4 Protonation:  not set
  5 Metals:       0 charge(s) set; cofactor edits for none
  6 Flexible:     none (rigid receptor)
  7 Rotatable:    not set
  8 Search box:   not set
  9 Settings:     not set

  r Run docking  (needs a receptor structure (1), a ligand (2), a search
    box (8), docking settings (9))
  c Docking campaign: run docking for every ligand of a library ...
  v View results: no run or campaign yet
  b Back

Enter choice (1):
\end{screen}

You work down the numbers. The \prompt{r} line always says what is still
needed, so you cannot forget a step.

\subsection{Step 1: the receptor}

The \term{receptor} is the protein the molecule is docked into.

\begin{steps}
  \item Type \key{1}. ProPrep names the loaded structure and asks
        \prompt{Use this structure as the receptor?} Type \key{y}.
  \item It lists the chains and asks which to include. HIV protease is a pair
        of identical chains and the pocket lies between them, so keep both:
        press \key{Enter} for all.
\begin{screen}
  Chains: A (99 residues), B (99 residues)
Chains to include (e.g. A or A,B; Enter for all) (all): §ans«C»
No chain C.
Chains to include (e.g. A or A,B; Enter for all) (all):
\end{screen}
\end{steps}

A chain letter that is not in the file is refused and the question comes
again.

\subsection{Step 2: the ligand}

The \term{ligand} is the small molecule to dock. There are three sources:
\key{s} a molecule already in the structure, \key{m} a \term{SMILES} string (a
molecule written as a line of text), or \key{f} a record of a molecule file
you loaded earlier.

\begin{steps}
  \item Type \key{2}, then \key{s}.
  \item ProPrep lists the candidates and you pick one:
\begin{screen}
Ligand source: [s] residue of the structure, [m] SMILES, [f] file (s): §ans«s»
┏━━━┳━━━━━━━━━┳━━━━━━┳━━━━━━━┓
┃ # ┃ residue ┃ name ┃ atoms ┃
┡━━━╇━━━━━━━━━╇━━━━━━╇━━━━━━━┩
│ 1 │ B:902   │ MK1  │ 45    │
└───┴─────────┴──────┴───────┘
Ligand number (range: 1-1, default: 1) (1): §ans«1»
\end{screen}
  \item ProPrep then describes the molecule it built, and where every detail
        came from:
\begin{screen}
  Ligand: 45 heavy atoms, 47 hydrogens, net charge +0; charged atoms: none
  SMILES: CC(C)(C)NC(=O)C1CN(Cc2cccnc2)CCN1CC(O)CC(Cc1ccccc1)C(=O)...
  - 45 heavy atoms from MK1 B902 in 1HSG.pdb.
  - Bond orders, formal charges and hydrogens from CCD MK1 (net charge +0).
\end{screen}
  \item Press \key{Enter} at the edit question to accept the molecule as it
        is.
\end{steps}

Choosing the ligand this way takes it out of the receptor automatically. The
dashboard confirms it: \prompt{MK1 B:902 is the ligand (2)}.

\begin{tip}
The edit question lets you change the chemistry by hand, for example
\typed{remove H9} to take off a hydrogen or \typed{charge O2 -1} to change a
charge. \typed{drop} takes back your last edit. You need this only for an
unusual protonation or redox state.
\end{tip}

\subsection{Step 4: protonation, and the chain-end question}
\label{sec:docking-protonation}

Whether a side chain carries a charge affects how the molecule binds. ProPrep
works this out with pdb2pqr and PROPKA at a pH you choose.

\begin{steps}
  \item Type \key{4} and press \key{Enter} to accept pH 7.0.
  \item ProPrep lists only the residues worth a second look, and says why:
\begin{screen}
  Running pdb2pqr with PROPKA at pH 7.0 ...
  pdb2pqr assigned 34 titratable residues.

  Of the 34 titratable residues, 31 are in their usual charged form at
  pH 7.0 (Asp-, Glu-, Lys+, ...) and 3 are not, or are His (which has no
  single usual form). Only those 3 are listed ('a' lists all). ...
┏━━━━━━━━━┳━━━━━━┳━━━━━━━━━━┳━━━━━━━━━━━━━━━━━━━━━━━━━━┓
┃ residue ┃ name ┃ template ┃ from                     ┃
┡━━━━━━━━━╇━━━━━━╇━━━━━━━━━━╇━━━━━━━━━━━━━━━━━━━━━━━━━━┩
│ A:69    │ HIS  │ HID      │ pdb2pqr/PROPKA at pH 7.0 │
│ B:25    │ ASP  │ ASH      │ pdb2pqr/PROPKA at pH 7.0 │
│ B:69    │ HIS  │ HID      │ pdb2pqr/PROPKA at pH 7.0 │
└─────────┴──────┴──────────┴──────────────────────────┘
\end{screen}
\end{steps}

The three-letter names are AMBER's forms of a residue: \texttt{HID} and
\texttt{HIE} are the two neutral forms of histidine (the hydrogen sits on one
ring nitrogen or the other), \texttt{HIP} is the positively charged form, and
\texttt{ASH} is aspartic acid with its hydrogen still on, so neutral. To
change one, type its identifier, for example \typed{A:69}, and then the name
you want.

\textbf{Now read the second table, the chain ends.} This is the one question
in the tool that a first run is likely to get wrong.

\begin{screen}
Chain ends: a real terminus is charged; a chain break is left neutral
┏━━━━━━━━━┳━━━━━━┳━━━━━━━┳━━━━━━━━━━━━━━━━━━━━━━━━━━━━━━━━━━━━━━━━━━┓
┃ residue ┃ name ┃ end   ┃ treated as                               ┃
┡━━━━━━━━━╇━━━━━━╇━━━━━━━╇━━━━━━━━━━━━━━━━━━━━━━━━━━━━━━━━━━━━━━━━━━┩
│ A:1     │ PRO  │ first │ N-terminus: charged NH3+ ...              │
│ A:99    │ PHE  │ last  │ C-terminus: charged COO- ...              │
│ B:1     │ PRO  │ first │ N-terminus: charged NH3+ ...              │
│ B:99    │ PHE  │ last  │ C-terminus: charged COO- ...              │
└─────────┴──────┴───────┴──────────────────────────────────────────┘
Residue or chain end to change (e.g. A:94), 'p' to click one in the
viewer, 'a' to list every titratable residue, or Enter when done ():
\end{screen}

A real chain end carries a charge. A \term{chain break} is where the chain
continues in the real protein but the atoms are missing from the file; such an
end is left neutral. A charged C-terminus needs an extra oxygen atom called
\term{OXT}. Many PDB files, \pdb{1HSG} among them, do not have it. If you
accept the proposal, the run fails.

\begin{watchout}
In version 1.23.0, ProPrep proposes \prompt{C-terminus: charged COO-} whenever
the sequence record says no residues are missing, even when the file has no
OXT atom. Docking then stops with a Meeko error that does not mention OXT:
\begin{screen}
No template matched for residue_key='A:99' ... CPHE heavy_miss=1 ...
Run docking failed:
Error: Creation of data structure for receptor failed.
Details:
- Template matching failed for: ['A:99', 'B:99']
\end{screen}
A Python error report is printed above it. The fix takes seconds and stays
inside ProPrep: open \key{4} again, type the residue named in the message
(\typed{A:99}), and answer \key{b}:
\begin{screen}
A:99 PHE, the last residue of chain A: [c] C-terminus/[b] chain break (c): §ans«b»
\end{screen}
Do the same for \typed{B:99}, press \key{Enter}, and run again. The other
route is to add the missing OXT atoms with the Structure Fixer
(Chapter~\ref{ch:cleaning}).
\end{watchout}

\subsection{Steps 3, 5, 6 and 7: what to leave alone}

\begin{itemize}
  \item \key{3} \textsf{Kept HETATM} decides which non-protein residues stay
        as part of the receptor. Because you picked MK1 as the ligand, there
        is nothing left to decide. Waters are left out by default.
  \item \key{5} \textsf{Metals \& cofactors} prints \prompt{No metals are kept
        in the receptor} for \pdb{1HSG} and has nothing to ask.
  \item \key{6} \textsf{Flexible} lets named side chains move while docking.
        Leave it at \prompt{none (rigid receptor)} for a first run.
  \item \key{7} \textsf{Rotatable} shows which bonds of the ligand are allowed
        to twist: 16 bonds for MK1, of which 14 rotate. The two amide bonds
        are held rigid, which is correct chemistry. Press \key{Enter} to
        accept; the dashboard then reads \prompt{7 Rotatable: reviewed}.
\end{itemize}

\subsection{Step 8: the search box}

The \term{search box} is the part of the protein ProPrep searches. A box
around the known pocket is both faster and more reliable than a box around the
whole protein. There are three ways to set it.

\begin{steps}
  \item Type \key{8}, then press \key{Enter} for \key{l} (around the ligand).
  \item Press \key{Enter} again to accept 8\,\AA{} of padding on every side.
\begin{screen}
Search box: [l] around the ligand, [r] around residues, [c] typed centre
and size (l):
Padding around them on every side (A) (min: 0.0, default: 8.0) (8.0):
  Box centre (11.63, 22.47, 5.86), edges 30.9 x 26.2 x 24.6 A (19925 A^3);
  the ligand as placed, padding 8 A.
\end{screen}
\end{steps}

The other two modes:

\begin{itemize}
  \item \key{r} builds the box around residues you name, which is what you
        want when the pocket is known but no molecule is in it. Typing
        \typed{A:25,B:25,A:50} (the two catalytic aspartates and one flap
        residue) gave a box of
        26.9\,$\times$\,30.1\,$\times$\,24.3\,\AA. A residue that is not in
        the structure is refused with \prompt{Not in the structure: A:999.}
        and the box stays unset.
  \item \key{c} asks for six numbers: the center and the three edge lengths.
\end{itemize}

\subsection{Step 9: the settings}

\begin{steps}
  \item Type \key{9}. ProPrep explains the three \term{scoring functions}
        (the formulas that turn a pose into a number) before asking:
\begin{screen}
  vina     AutoDock Vina's empirical function (Trott & Olson, 2010): steric,
           hydrophobic and hydrogen-bond terms. No partial charges and no
           electrostatics. Fast.
  vinardo  A re-parameterisation of Vina's terms (Quiroga & Villarreal,
           2016); same speed.
  ad4      AutoDock 4's force-field function (Huey et al., 2007): van der
           Waals, hydrogen bonds, desolvation and electrostatics from the
           Gasteiger charges, so charges and metal ions count. Needs
           autogrid4 maps, built once per receptor and box, so it is slower.
  Scores are in kcal/mol for all three but are not comparable between
  functions.
Scoring function: vina, vinardo or ad4 (vina): §ans«vina»
\end{screen}
  \item Press \key{Enter} through the rest for a first run, or set them as we
        did for the fast runs: seed \typed{7}, \term{exhaustiveness}
        \typed{4}, \typed{5} poses.
\end{steps}

\begin{tabularx}{\linewidth}{@{}lL@{}}
\toprule
\textbf{Setting (default)} & \textbf{What it means} \\
\midrule
Scoring function (vina) & Which formula scores a pose. Start with
  \texttt{vina}. \\
Random seed (42) & Docking uses random numbers; the same seed repeats the
  same run. Not 0, which Vina reads as ``pick one''. \\
Exhaustiveness (8) & How hard to search. Higher is slower and more thorough.
  4--8 was enough here. \\
Number of poses (9) & How many positions to report. \\
Energy range (3.0) & Only poses within this many kcal/mol of the best are
  kept. \\
Minimum RMSD between poses (1.0) & How different two reported poses must be. \\
Grid spacing (0.375) & Only used by \texttt{ad4}. \\
CPU cores (0) & 0 means use all of them. \\
\bottomrule
\end{tabularx}

\subsection{Run it}

Type \key{r}. ProPrep prepares both sides and runs the docking.

\begin{screen}
  Building the ligand and receptor and docking; this takes seconds to
  minutes ...
  Done in 17.4 s. Output: .../docking/docking_20261121_131645
  - Vina: WARNING: At low exhaustiveness, it may be impossible to utilize
    all CPUs.
  - 5 poses within 3.0 kcal/mol of the best, vina scoring, seed 7,
    exhaustiveness 4.
Docked poses (best first)
┏━━━━━━┳━━━━━━━━━━━━━━━━━━┳━━━━━━━━┳━━━━━━━┳━━━━━━━━━━┳━━━━━━━━━━━━━━━━━━━━━┓
┃ pose ┃ total (kcal/mol) ┃ inter  ┃ intra ┃ torsions ┃ RMSD to crystal (A) ┃
┡━━━━━━╇━━━━━━━━━━━━━━━━━━╇━━━━━━━━╇━━━━━━━╇━━━━━━━━━━╇━━━━━━━━━━━━━━━━━━━━━┩
│ 1    │ -11.56           │ -20.35 │ -1.78 │ 8.79     │ 0.46                │
│ 2    │ -10.67           │ -18.95 │ -1.60 │ 8.11     │ 11.07               │
│ 3    │ -10.59           │ -18.44 │ -1.97 │ 8.05     │ 4.16                │
│ 4    │ -10.49           │ -17.94 │ -2.30 │ 7.97     │ 10.35               │
│ 5    │ -10.12           │ -17.88 │ -1.71 │ 7.69     │ 10.52               │
└──────┴──────────────────┴────────┴───────┴──────────┴─────────────────────┘
  The viewer is not open; 'v' shows any pose in it.
\end{screen}

\section{Understanding the results}
\label{sec:docking-results}

\paragraph{The columns.} \textbf{total} is the score, in kcal/mol; more
negative means a better predicted fit. \textbf{inter} is the part from contact
between the ligand and the protein, \textbf{intra} the part from the ligand's
own shape, and \textbf{torsions} a penalty for every bond that had to twist.
\textbf{RMSD to crystal} appears only when the ligand came from the structure:
it is how far this pose is from where the crystal structure put it.

\paragraph{The best pose is not always first.} We ran the same system three
times, changing only the scoring function:

\begin{tabularx}{\linewidth}{@{}llllL@{}}
\toprule
\textbf{Scoring} & \textbf{Time} & \textbf{Best score} & \textbf{RMSD} &
\textbf{What it means} \\
\midrule
vinardo & 29.1\,s & $-9.07$ & 0.41\,\AA{} & Top pose is the crystal pose. \\
vina & 17.4\,s & $-11.56$ & 0.46\,\AA{} & Top pose is the crystal pose. \\
ad4 & 59.0\,s & $-12.88$ & 4.05\,\AA{} & Top pose is \emph{wrong}; pose 5
  ($-11.24$) was the right one, at 1.78\,\AA. \\
\bottomrule
\end{tabularx}

\medskip
Two lessons. First, scores from different functions cannot be compared: $-9.07$
from vinardo and $-11.56$ from vina describe the same pose. Second, the
top-ranked pose is not automatically the answer. Look at the first few poses,
and if you have a crystal pose to compare with, let the RMSD column decide.

\paragraph{The files.} Every run writes its own folder, named with the date and
time, inside \file{docking/}:

\begin{tabularx}{\linewidth}{@{}lL@{}}
\toprule
\textbf{File} & \textbf{What it holds} \\
\midrule
\file{poses.sdf} & Every pose, best first, each carrying its rank, energy and
  RMSD. \\
\file{pose_1.sdf} \dots & One file per pose, for a viewer or another program. \\
\file{docking.json} & The complete setup, including every protonation
  decision. Your record of what you did. \\
\file{vina_output.txt} & Vina's own output, as printed. \\
\file{ligand.pdbqt}, \file{receptor_rigid.pdbqt} & The prepared molecules, for
  anyone who wants to re-run Vina by hand. \\
\bottomrule
\end{tabularx}

\medskip
An \texttt{ad4} run also writes the AutoGrid map files. A run with flexible
side chains adds a second receptor file,
\file{receptor_flex.pdbqt}, and one file per pose holding the side chains as
they were docked.

\paragraph{Looking at a pose.} Type \key{v}, then the pose number, to open it
in the browser viewer with the box drawn around it. Answer \typed{0} to see
only the table.

\section{Other choices}

\subsection{Letting side chains move}

Pockets are not rigid. Option \key{6} lets you name side chains that may
change shape while docking. We tried the catalytic Asp25 and Ile50 of the flap:

\begin{screen}
Residues whose side chains move during docking (e.g. A:45,A:88; 'none' for
a rigid receptor; 'keep' to leave as is; ...) (keep): §ans«A:25,A:50»
\end{screen}

The run took 26.9\,s and gave a best pose of $-11.43$ at 0.50\,\AA{} RMSD:
about the same as the rigid run here. Flexible side chains cost time and help
mainly when you expect the pocket to adjust.

\subsection{A ligand from SMILES}

Type \key{2}, then \key{m}, then the SMILES string. ProPrep builds a
three-dimensional shape with RDKit and tidies it up:

\begin{screen}
SMILES: §ans«CC(=O)Nc1ccc(O)cc1»
Seed for the 3D build (RDKit ETKDG) (min: 1, default: 42) (42):
Minimise the 3D build with MMFF94? [mmff94/none] (mmff94):

  Ligand: 11 heavy atoms, 9 hydrogens, net charge +0; charged atoms: none
  - 3D coordinates from RDKit ETKDGv3, seed 42.
  - Atom names generated for all 20 atoms (element + number).
  - Geometry minimised with MMFF94 (converged).
\end{screen}

\begin{watchout}
A molecule built from SMILES has no position in your protein: it sits
wherever RDKit put it, near the origin of the coordinates. Box mode \key{l}
(around the ligand) then puts the box in empty space. ProPrep warns you:
\prompt{A ligand built from SMILES sits wherever RDKit put it, not in the
receptor; a box around it is rarely what you want.} If you run it anyway, the
docking \emph{succeeds} with no error, no RMSD column and poses of about
$-4.15$~kcal/mol. The only hints that nothing was docked are the time
(0.7~seconds) and the weak scores. For a SMILES ligand always use \key{r}
(around named residues) or \key{c} (typed numbers).
\end{watchout}

\begin{watchout}
Changing the ligand clears two earlier settings: the box (option \key{8}),
because it belonged to the old ligand, and the keep/leave-out list (option
\key{3}). Set the box again after every ligand change.
\end{watchout}

\begin{watchout}
If you leave the crystal ligand kept in the receptor (option \key{3}) and dock
something else into the same pocket, the site is already occupied and your
scores are meaningless. Choosing a ligand with \key{2} \key{s} removes it for
you; a ligand from SMILES or a file does not.
\end{watchout}

\subsection{Many ligands at once: a campaign}

A \term{campaign} runs the same receptor, box and settings against every
molecule in a library file.

\begin{steps}
  \item Load the library first: \key{1a}, then \key{4} (\textsf{Load a ligand
        file or library (SDF/mol2/SMILES) for docking}), then type the full
        path to the file.
\begin{screen}
  mini_library.smi: 3 molecule(s)
    0: paracetamol
    1: benzoic_acid
    2: not_a_real_drug
Stored for Molecular Docking: .../mini_library.smi
\end{screen}
  \item Back in \key{5d}, set the receptor, the box and the settings as before
        (a campaign needs no single ligand), then type \key{c}.
  \item Answer the questions about the library as a whole: what to do with
        molecules whose three-dimensional shape is not given, which seed to
        use, whether to minimize, and whether rings may flex. The defaults
        were right for our library.
  \item At \prompt{Start docking? (Ctrl-C stops; finished ligands are kept)}
        type \key{y}.
\end{steps}

\begin{screen}
  3 molecules in mini_library.smi; 0 already done.
  Prepared 3; skipped 0 (reasons in results.csv).
  [1/3] paracetamol: -5.34 kcal/mol (1 s)
  [2/3] benzoic_acid: -5.16 kcal/mol (1 s)
  [3/3] not_a_real_drug: -4.70 kcal/mol (1 s)
┏━━━━━━┳━━━━━━━━━━━━━━━━━┳━━━━━━━━━━━━━━━━━━┳━━━━━━━━━━━━━━━━┳━━━━━━━━━━━━━┓
┃ rank ┃ ligand          ┃ score (kcal/mol) ┃ per heavy atom ┃ heavy atoms ┃
┡━━━━━━╇━━━━━━━━━━━━━━━━━╇━━━━━━━━━━━━━━━━━━╇━━━━━━━━━━━━━━━━╇━━━━━━━━━━━━━┩
│ 1    │ paracetamol     │ -5.339           │ -0.485         │ 11          │
│ 2    │ benzoic_acid    │ -5.155           │ -0.573         │ 9           │
│ 3    │ not_a_real_drug │ -4.696           │ -0.587         │ 8           │
└──────┴─────────────────┴──────────────────┴────────────────┴─────────────┘
\end{screen}

The campaign folder holds \file{results.csv} (one row per molecule, with the
score, the score per heavy atom, the number of poses and the reason for any
skip), \file{top_poses.sdf} (each molecule's best pose, best first), and a
folder per molecule with its own poses. \file{manifest.json} records what the
campaign was made from, so an interrupted campaign can be continued if nothing
has changed.

\begin{tip}
The \prompt{per heavy atom} column divides the score by the number of heavy
atoms. Bigger molecules score better for being bigger alone; this column is a
fairer comparison between molecules of different sizes.
\end{tip}

Type \key{v} to see the results again. When both a single run and a campaign
exist, ProPrep asks which you mean. Choosing a rank from a campaign writes
\file{best_pose.sdf} next to that molecule's poses.

\section{If something goes wrong}

\begin{tabularx}{\linewidth}{@{}lL@{}}
\toprule
\textbf{What you see} & \textbf{What to do} \\
\midrule
\prompt{Needs a loaded structure} & Load a structure with \key{1a} first. \\
\prompt{Template matching failed for: ['A:99', 'B:99']} & The C-terminus has
  no OXT atom. Set those ends to chain break in option \key{4}
  (Section~\ref{sec:docking-protonation}). \\
\prompt{No chain C.} / \prompt{Not in the structure: A:999.} & A chain letter
  or residue that is not in the file. Check the lists ProPrep printed. \\
Docking finishes in under a second with weak scores & The box probably does
  not contain the protein. Set it again with \key{r} or \key{c}. \\
\prompt{The ligand could not be built: Unsupported ligand file type '.smi'} &
  A \file{.smi} library works for a campaign but not as a single ligand. Use
  \file{.sdf}, \file{.mol} or \file{.mol2}, or run a campaign. \\
Residues reported as missing heavy atoms & Rebuild them with the Structure
  Fixer first (Chapter~\ref{ch:cleaning}). \\
\bottomrule
\end{tabularx}

\begin{watchout}
In version 1.23.0, a refused ligand leaves its name on the dashboard. After
the \file{.smi} message above, line \key{2} still read \prompt{Ligand:
mini\_library.smi record 1} although no ligand had been built. Check the
dashboard before you run.
\end{watchout}
